\documentclass[]{article}
\usepackage{longtable}

\input{preamble.tex}

\colorlet{accent}{jimp}

\title{Dokumentacja implementacyjna - Graficzny interfejs użytkownika dla wizualizacji grafów w języku Java}
\author{\authoricon Aleksander Jóźwik, Anastasiya Kryvetskaya}
\date{\dateicon \today}

\definecolor{zalety}{HTML}{a6d189}
\definecolor{wady}{HTML}{e78284}

\begin{document}
	\noindent
	\maketitle

	\tableofcontents
	\clearpage

    \section{Klasa  Main}
    Odpowiada za uruchomienie UI programu.
    \section{Klasa GraphVisualizerUI}
    Odpowiada za stworzenie UI korzystając z innych klas. Posiada dwie funkcje prywatne:
    \begin{itemize}
        \item initFrame - ustawia właściwości okna tj. rozmiar, tytuł, układ rozmieszczenia elementów czy ikonę.
        \item initLayout - dodaje pasek menu, lewe i prawe okno oraz pasek stanu.
    \end{itemize}
    \section{Klasa AppTheme}
    Odpowiada za zdefiniowanie palety kolorów i czcionek.
    \section{Klasa TopMenuBar}
    Odpowiada za górny pasek menu. Posiada dwie funkcje prywatne:
    \begin{itemize}
        \item createStyledMenu - tworzy nowy menu zgodnie z paletą kolorów.
        \item createStyledMenuItem - tworzy nowy element menu zgodnie z paletą kolorów.
    \end{itemize}
    \section{Klasa SidebarPanel}
    Odpowiada za boczny panel pozwalający na konfiguracje parametrów programu i uruchomienie wizualizacji.
    \begin{itemize}
        \item buildUI - tworzy elementy UI w bocznym panelu.
        \item createColorSwatch - tworzy przycisk odpowiadający za zmianę koloru.
        \item addPlainText - służy do tworzenia zwykłego tekstu.
        \item addSectionText - służy do tworzenia tekstu nagłówkowego.
        \item createRowWithLabel - tworzy panel, w którym umieszcza etykietę i element podany jako argument.
    \end{itemize}
    \section{Klasa StatusBar}
    Odpowiada za dolny pasek stanu.
    \section{Klasa GraphCanvas}
    Odpowiada za płótno na którym wyświetalana jest wizualizacja.
    \section{Klasa AboutWindow}
    Odpowiada za okno "O Programie".
\end{document}